\subsection{Jets of maps of Banach spaces}

In this section, E and F denote two Banach spaces. If m is an integer \(\geqslant0\) , one denotes by \(P_{m}(E;F)\) the Banach space of continuous homogeneous polynomials of degree m on E with values in F (Part 1, Appendix, A.2).

12.2.1. Let U be an open subset of E, let \(f\colon\mathbf{U}\to\mathbf{F}\) be a map of class \(\mathbf{C}^{r}\), and let \(a\in\mathbf{U}\). There exists one and only one continuous polynomial

\[
\tilde {f} = f _ {0} + \dots + f _ {k} \quad (\text{with }f _ {m} \in \mathrm{P} _ {m} (\mathrm{E}; \mathrm{F}))
\]

of degree \(\leqslant k\), having the same jet of order k at the origin as \(x \mapsto f(x - a)\). If \(0 \leqslant m \leqslant k\), the \(m\)-th component \(f_m\) of \(\tilde{f}\) is denoted \(\Delta^m f(a)\) or \(\Delta_a^m(f)\). When \(r = \omega\), this notation coincides with that of 3.2.1 and 4.2.1; when \(r \leqslant \infty\), one has

\[
m! \Delta^ {m} f (a) (h) = \mathrm{D} ^ {m} f (a) (h, \dots , h) = \mathrm{D} ^ {m} f (a). h ^ {m}.
\]

The map \(f\mapsto\tilde{f}\) defines, by passing to the quotient, a bijection from \(\mathrm{J}_{a}^{k}(\mathrm{E},\mathrm{F})\) onto \(\prod_{0\leqslant m\leqslant k}\mathrm{P}_{m}(\mathrm{E};\mathrm{F})\) by means of which one identifies these two spaces; in particular, \(\mathrm{J}_{a}^{k}(\mathrm{E},\mathrm{F})\) is equipped with a Banach space structure over K. If \(j\in\mathrm{J}_{a}^{k}(\mathrm{E},\mathrm{F})\) , one denotes by \(\Delta_{a}^{m}(j)\) the \(m\) -th component of \(j\) ; we have

\[
\Delta_ {a} ^ {m} (j) \in \mathrm{P} _ {m} (\mathrm{E}; \mathrm{F}) \quad \text {for} \quad 0 \leqslant m \leqslant k, \quad \text {and} \quad \Delta_ {a} ^ {0} (j) = b (j) \in \mathrm{F}.
\]

12.2.2. Let U (resp. V) be an open subset of E (resp. F). Denote by \(Q^{k}(E, F)\) the product Banach space of the \(P_{m}(E; F)\) for \(1 \leqslant m \leqslant k\). The map

\[
j \mapsto (s (j), b (j), \Delta_ {s (j)} ^ {1} (j), \dots , \Delta_ {s (j)} ^ {k} (j))
\]

is a bijection of \(\mathbf{J}^{k}(\mathbf{U},\mathbf{V})\) onto \(\mathbf{U}\times\mathbf{V}\times\mathbf{Q}^{k}(\mathbf{E},\mathbf{F})\), by means of which these two sets are identified. In particular, \(\mathbf{J}_{0}^{k}(\mathbf{E},\mathbf{F})_{0}\) is identified with \(\mathbf{Q}^{k}(\mathbf{E},\mathbf{F})\).

